\documentclass[12pt]{scrartcl}

%Matthis Vorlage
\usepackage{../mathe_definition}

%Ausgelagerter Code fürs schöne seitenlayout
\date{\today}
\title{{\itshape Stochastische Prozesse}}
\module{Blatt 1}
\author{Philip Wunderlich \and Jonathan Langer}
%\matriculation{3737927\and 3750030}

\fancypagestyle{netter_header}{
	\renewcommand{\and}{\\}
	\fancyhead[L]{\MyAuthor}
	\fancyhead[C]{\MyTitle\\ \MyModule}
	\fancyhead[R]{\MyDate\\Seite \thepage{} von \pageref{LastPage}}
	\fancyfoot[C]{}
}
\pagestyle{netter_header}


%Abkürzungen
\newcommand{\einruecken}{\begin{addmargin}[25pt]{0pt}} 
	%LaTeX mag keine Umlaute :(
\newcommand{\ausruecken}{\end{addmargin}}
\newcommand{\R}{\ensuremath{\mathbb{R}}}
\newcommand{\Om}{\ensuremath{\Omega}}
\newcommand{\A}{\ensuremath{\mathcal{A}}}
\newcommand{\B}{\ensuremath{\mathcal{B}}}
\newcommand{\C}{\ensuremath{\mathcal{C}}}
\newcommand{\WraumAnn}{\ensuremath{(\Om, \A, P)} ein Wahrscheinlichkeitsraum}
\newcommand{\MraumAnn}{\ensuremath{(E,\B)} ein Messraum}
\newcommand{\Xt}{\ensuremath{(X_t)_{t\in T}}}
\newcommand{\Yt}{\ensuremath{(Y_t)_{t\in T}}}

\begin{document}

\section*{Aufgabe 1}
\subsubsection*{Lemma 1.1.3}
\einruecken
	
Seien \WraumAnn{} und \MraumAnn{}.\\
Ferner sei \Xt{} ein stochastischer Prozess.\\
Dann ist $X_{\bullet}:\Om\rightarrow E^T$ $(\A, \bigotimes_{t\in T}\B)$-messbar
	
\ausruecken

\subsection*{Beweis}
\einruecken
	
per Definition gilt $\bigotimes_{t\in T}\B=\sigma(\mathcal{Z})$ mit $\mathcal{Z}$ die Zylindermengen.\\
Sei $A \in \mathcal{Z}$. Dann $A=\prod_{t\in T}A_t$ mit $A_t \in \B$ und endlichem $J \subset T$, sodass $A_t=\Omega_t \forall t\in T\setminus J$.\\
Sei $B_{t,s} :=
\begin{cases*}
	A_t & falls t=s \\
	\Omega & sonst
\end{cases*}$.\\
Dann $$X_{\bullet}^{-1}(\prod_{t\in T}B_{t,s})=X_s^{-1}(A_s) \text{ für } s\in T$$
Wegen $$A=\prod_{t\in T}A_t=\bigcap_{s\in T}\prod_{t\in T}B_{t,s}$$
gilt also 
\begin{align*}
	X_{\bullet}^{-1}(A) &= X_{\bullet}^{-1}(\bigcap_{s\in T}\prod_{t\in T}B_{t,s})\\
	&= \bigcap_{s\in T} X_{\bullet}^{-1}(\prod_{t\in T}B_{t,s})\\
	&= \bigcap_{s\in T} X_s^{-1}(A_s)
\end{align*}
Also $X_{\bullet}$ messbar, da $X_t$ messbar für alle $t \in T$.
	
\ausruecken

\section*{Aufgabe 1}
\subsubsection*{Lemma 1.1.7}
\einruecken

Seien \WraumAnn{} und \MraumAnn{}. Ferner seien $X := \Xt$ und $Y := \Yt$ zwei E-wertige, stochastische Prozesse über $\Omega$. Dann gilt:\\

iii) Für T = [0,1] gilt die Implikation
\einruecken
\textit{X Version von Y $\Rightarrow$ X und Y ununterscheidbar}
\ausruecken
nicht.

\ausruecken

\subsection*{Beweis}
\einruecken

trivial

\ausruecken

\section*{Aufgabe 1}
\subsubsection*{Lemma 1.1.9}
\einruecken

Sei \WraumAnn{} und $X := \Xt$ ein stetiger, stochastischer Prozess. Dann gibt es einen von X ununterscheidbaren stochastischen Prozess $Y := \Yt$, für den jeder Pfad stetig ist.

\ausruecken

\subsection*{Beweis}
\einruecken

trivial

\ausruecken

\subsubsection*{Satz 1.1.10}
\einruecken

Seien $X := \Xt$ und $Y := \Yt$ (rechts-/linksseitig) stetige, stochastische Prozesse über einem Wahrscheinlichkeitsraum $(\Om,\A,P)$.\\
Ist Y eine Version von X, so sind X und Y ununterscheidbar.

\ausruecken

\subsection*{Beweis}
\einruecken

trivial

\ausruecken

\section*{Aufgabe 1}
\subsubsection*{Satz 1.1.13}
\einruecken

Sei \WraumAnn{} und $X := \Xt$ ein \R-wertiger stochastischer Prozess über $\Omega$.\\
Dann gelten die folgenden Aussagen:\\

i) Für jedes $n\ge 1$, alle Permutationen $\pi:{1,\dots,n}\rightarrow{1,\dots,n}$ und alle paarweise verschiedenen $t_1,\dots,t_n\in T$ gilt
$$P_{\pi(t_1,\dots,t_n)}(A) = P_{(t_1,\dots,t_n)}(\pi^{-1}(A)),
\qquad A\in \B^n.$$

ii) Für alle $n\ge m\ge 1$ und alle paarweise verschiedenen $t_1,\dots,t_n\in T$ gilt
$$P_{(t_1,\dots,t_m)}(A) = P_{(t_1,\dots,t_n)}(A\times \R^{n-m}),
\qquad A \in \B^m.$$

\ausruecken

\subsection*{Beweis}
\einruecken

trivial

\ausruecken

\section*{Aufgabe 5}
\einruecken

Seien \WraumAnn{} und $\B,\C \subset \A$ $\sigma$-Algebren.\\
Man nehme an, \B{} sei \textit{P-trivial}, i.e. $P(C)\in {0,1}$ für alle $C\in\C$.\\
Zeigen Sie dass \B{} und \C{} unabhängig sind.


\ausruecken
\subsection*{Beweis}
\einruecken

tbd

\ausruecken

\end{document}